\section{Introduction}

\subsection{General Current Status and Public Health Motivation (G)}
With the rapid expansion of Artificial Intelligence of Things (AIoT) sensing infrastructure and open government environmental repositories, public access to ambient air quality monitoring has improved substantially worldwide. Nevertheless, outdoor air pollution persists as a critical public health threat. According to the World Health Organization (WHO), approximately 99\% of the global population in 2019 lived in areas exceeding WHO Air Quality Guideline (AQG) limits, resulting in roughly 4.2 million premature deaths annually \cite{who2021,who2024}. In Taiwan, long-term air quality management documented by the Ministry of Environment (MOENV) reduced the national annual mean concentration of fine particulate matter ($\text{PM}_{2.5}$) from $16.2\,\mu\text{g/m}^3$ in 2019 to $12.8\,\mu\text{g/m}^3$ in 2025, representing an approximate 21\% reduction \cite{moenv2026news,moenv2026stat}. Taiwan's MOENV Open Data Platform continuously disseminates station-level real-time Air Quality Index (AQI) reports (\texttt{AQX\_P\_432}), daily updated hourly pollutant series (\texttt{AQX\_P\_13}), historical AQI and wind archives (\texttt{AQX\_P\_488}), and TWD97 station coordinates (\texttt{AQX\_P\_07}) \cite{aqxp432,aqxp488,aqxp13,aqxp07}, creating a rich empirical foundation for data-driven environmental services without requiring proprietary sensor hardware.

\subsection{However: Core Technical and Domain Challenges (H)}
However, an encouraging national annual average operates on a fundamentally different spatiotemporal scale than an individual's immediate daily activity decision, such as selecting a 40-minute window over the next three to six hours for outdoor jogging or commuting near a specific monitoring station in Taichung City. Translating raw monitoring data into reliable personal activity guidance faces four severe challenges:
\begin{enumerate}
    \item \emph{Spatiotemporal and Timestamp Provenance Ambiguity:} Existing platforms routinely conflate observation time (\texttt{observed\_at}), official release time (\texttt{published\_at}), and client fetch time (\texttt{retrieved\_at}). Upstream sensor maintenance or telemetry gaps produce missing hourly slots that naive pipelines erroneously fill with $0\,\mu\text{g/m}^3$, falsely indicating pollution-free air. Furthermore, point-source regulatory stations characterize neighborhood-scale ambient concentrations and cannot independently substantiate street-by-street clean-route rankings without co-located mobile sensors \cite{steinle2018,epa2024}.
    \item \emph{Regulatory Divergence Between Official AQI and Single-Pollutant $\text{PM}_{2.5}$:} Under Taiwan's MOENV AQI framework (including the revised sub-index concentration breakpoints effective January 1, 2025), official AQI is the maximum sub-index across six pollutants ($\text{O}_3$, $\text{PM}_{2.5}$, $\text{PM}_{10}$, $\text{CO}$, $\text{SO}_2$, and $\text{NO}_2$) computed over distinct moving-average windows \cite{moenv2025aqi}. Equating a single-pollutant raw hourly $\text{PM}_{2.5}$ forecast to a ``predicted official AQI'' violates regulatory semantics.
    \item \emph{Time-Series Leakage and the Accuracy--Decision Gap:} Many $\text{PM}_{2.5}$ forecasting pipelines employ random train--test splitting or future-looking rolling features that leak future distribution statistics, while omitting simple baselines such as Persistence and Seasonal Naive \cite{chianese2026,hewamalage2023,lu2026}. Moreover, empirical evidence shows that lower overall Mean Absolute Error (MAE) does not automatically guarantee better outdoor scheduling decisions when candidate windows differ by less than model uncertainty or when acute pollution spikes are smoothed \cite{berlinghieri2026}.
    \item \emph{Unconstrained Generative AI Risks:} Deploying an unconstrained LLM chatbot over environmental APIs invites numeric hallucination, silent substitution of stale cache data as live observations, vulnerability to indirect prompt injection embedded in external API payloads, and unauthorized medical advice \cite{yao2023,nie2025,owasp2026}.
\end{enumerate}

\subsection{In Literature and Judgement of Existing Paradigms (I, J)}
In the literature, approaches to environmental decision support have evolved from simple to complex across three major categories: 1) static open-data dashboards and deterministic rule tables \cite{aqxp432,moenv2025aqi}; 2) statistical, classical ML, and deep spatiotemporal forecasting models paired with post-hoc explainable AI (XAI) such as SHAP \cite{chianese2026,joseph2025,lundberg2017,houdou2024,lo2025}; and 3) tool-augmented LLM agents utilizing Reasoning and Acting (ReAct) or Instructor--Worker workflows \cite{yao2023,nie2025,gao2025}.

Evaluating the strengths, weaknesses, opportunities, and threats (SWOT Judgement) of these paradigms reveals a clear trade-off:
\begin{itemize}
    \item \emph{Rule-Based Dashboards:} Offer high regulatory fidelity and zero hallucination (Strength), but provide only backward-looking or current status without multi-horizon candidate window comparison (Weakness).
    \item \emph{Standalone ML/DL Forecasters:} Capture nonlinear temporal autocorrelation across 1--6 hour horizons (Strength), yet frequently lack strict rolling-origin baseline audits, conflate SHAP feature attribution with physical atmospheric causality, and fail to translate point errors into uncertainty-aware activity abstention (Weakness/Threat) \cite{hewamalage2023,houdou2024,berlinghieri2026}.
    \item \emph{Autonomous LLM Agents:} Excel at natural language intent parsing and multi-step tool orchestration (Opportunity), but remain brittle under missing/stale data and vulnerable to prompt injection or regulatory overclaiming when not bounded by deterministic domain guardrails (Threat) \cite{nie2025,owasp2026}.
\end{itemize}

\subsection{Keypoints, Novelty, and Paper Organization (K, L)}
To bridge this gap, the core design spirit and novelty of AirGuard AI (v3 specification) lie in a \emph{three-tier separation of concerns}: an allowlist-constrained ReAct planner agent handles intent interpretation and tool routing, version-locked leakage-safe ML models generate 1-, 3-, and 6-hour raw hourly $\text{PM}_{2.5}$ forecasts, and a deterministic domain rule engine enforces 2025 MOENV AQI activity guidance, concentration-time exposure indexing, and fail-closed safety degradation. The major contributions of this paper are summarized as follows:
\begin{enumerate}
    \item \emph{Reproducible MOENV Data Engineering and Cryptographic Provenance (Verified in Phase W1):} We implement and empirically audit a Python ingestion and quality-control pipeline across four MOENV datasets for all six official stations in Taichung City (123,055 valid hourly observations, 2024--2026), resolving duplicate daily rows, trailing date whitespace, and unobserved calendar days while saving raw paginated bytes with SHA256 manifests for deterministic offline replay.
    \item \emph{Leakage-Safe Multi-Horizon $\text{PM}_{2.5}$ Benchmarking Protocol:} We design a chronological 70/15/15 split and rolling-origin evaluation comparing Persistence and 24-hour Seasonal Naive baselines against Ridge regression, Random Forest, and XGBoost at $h \in \{1, 3, 6\}\text{ h}$ on identical valid sample pairs, tagging historical backtests without publication-delay logs as \texttt{reconstructed\_from\_observed\_at}.
    \item \emph{Dual-Track Evidence Separation and Conservative Window Comparison:} We architect an eight-module pipeline that strictly separates official multi-pollutant AQI from predicted hourly $\text{PM}_{2.5}$ and station-scale concentration-time exposure proxies ($\mu\text{g}\cdot\text{h/m}^3$), abstaining from forced window rankings when inter-window differences fall within validation error bounds.
    \item \emph{Dual-Track Evaluation Matrix and Scope Governance:} We couple quantitative forecasting metrics (MAE, RMSE, Skill Score, decision regret) with a 30-scenario agent verification suite testing tool accuracy, citation completeness, regulatory compliance, prompt-injection defense, and fail-closed degradation.
\end{enumerate}

The remainder of this paper is organized as follows. Section~II reviews related work across AI technology and environmental domain know-how, synthesizing our Deep Research findings and structured comparison table. Section~III details the proposed AirGuard AI design following the systematic \texttt{PA}--\texttt{PM} methodology. Section~IV presents the Phase W1 empirical data audit and the dual-track experimental evaluation framework. Section~V concludes the paper and outlines ongoing work.
